%% LyX 2.3.6.2 created this file.  For more info, see http://www.lyx.org/.
%% Do not edit unless you really know what you are doing.
\documentclass[english]{article}
\usepackage[T1]{fontenc}
\usepackage[latin9]{inputenc}
\usepackage{amsmath}
\usepackage{amssymb}
\usepackage{babel}
\begin{document}
As detailed in the main text, the starting point of the resistivity
calculation is the variational Boltzmann approach justified by the
long mean free path $k_{F}\ell$. Within this approach the transport
lifetime can estimated by 
\begin{equation}
\rho\leq\frac{\frac{V}{2kT}\sum_{\boldsymbol{q}\boldsymbol{p}}P_{\boldsymbol{p}\boldsymbol{q}}\left(\psi_{\boldsymbol{q}}-\psi_{\boldsymbol{p}}\right)^{2}}{\left|\sum_{\boldsymbol{q}}e\frac{\partial f}{\partial\epsilon_{\boldsymbol{q}}}\psi_{\boldsymbol{q}}\boldsymbol{v}\left(\boldsymbol{q}\right)\cdot\boldsymbol{n}\right|^{2}}
\end{equation}

with $\psi_{\boldsymbol{q}}$ a variational anzatz, $P_{\boldsymbol{p}\boldsymbol{q}}$
the transition probaility from different momentum states,$V$ the
volume, $\boldsymbol{n}$ is the direction of the electric field $\boldsymbol{E}$,
$f$ the fermi-dirac distribution, $e$ the electron charge , $\boldsymbol{v}\left(\boldsymbol{q}\right)=\partial_{\boldsymbol{q}}\epsilon_{\boldsymbol{q}},$
and $\epsilon_{\boldsymbol{q}}$ the dispersion. Using the ansatz
$\psi_{\boldsymbol{q}}=-e\boldsymbol{v}_{\boldsymbol{q}}\cdot\boldsymbol{E}\tau_{\text{tr}}$
we recover the drude formula with the momentum relaxation time
\[
\tau_{\text{tr}}^{-1}=\frac{\frac{1}{kT}\sum_{\boldsymbol{q}\boldsymbol{p}}P_{\boldsymbol{p}\boldsymbol{q}}\left(\boldsymbol{v}\left(\boldsymbol{q}\right)\cdot\boldsymbol{n}-\boldsymbol{v}\left(\boldsymbol{q}\right)\cdot\boldsymbol{n}\right)^{2}}{\sum_{\boldsymbol{q}}2\left(-\frac{\partial f}{\partial\epsilon_{\boldsymbol{q}}}\right)\left(\boldsymbol{v}\left(\boldsymbol{q}\right)\cdot\boldsymbol{n}\right)^{2}}.
\]

Assuming Mathiessen's rule holds, we have 
\begin{equation}
\tau_{\text{tr}}^{-1}=\tau_{\text{tr},K}^{-1}+\tau_{\text{tr},EME}^{-1}+\tau_{\text{tr},ep}^{-1}
\end{equation}

Where each individual scattering rate is determined by the transition
probabilities 
\begin{equation}
P_{\boldsymbol{p}\boldsymbol{q}}^{K}=kT\frac{2\pi}{\hbar}\left(-\frac{\partial f_{\boldsymbol{p}}}{\partial\epsilon_{\boldsymbol{p}}}\right)\left|J_{K}\left(\boldsymbol{p},\boldsymbol{q}\right)\right|^{2}\chi''\left(\boldsymbol{p}-\boldsymbol{q},\varepsilon_{\boldsymbol{p}}-\epsilon_{\boldsymbol{q}}\right)\left(n_{B}\left(\varepsilon_{\boldsymbol{p}}-\epsilon_{\boldsymbol{q}}\right)+n_{F}\left(-\epsilon_{\boldsymbol{q}}\right)\right)
\end{equation}

\begin{equation}
P_{\boldsymbol{k}\boldsymbol{p}}^{EME}=kT\left(-\frac{\partial f}{\partial\epsilon_{\boldsymbol{p}}}\right)\left|\tilde{\alpha}\left(\boldsymbol{p},\boldsymbol{q}\right)\right|^{2}\sum_{\boldsymbol{k}}\int\frac{d\nu}{\pi}\chi''\left(\boldsymbol{k}+\boldsymbol{p}-\boldsymbol{q},\epsilon_{\boldsymbol{p}}-\nu\right)\mathcal{B}\left(\boldsymbol{k},\nu-\varepsilon_{\boldsymbol{q}}\right)\left[n_{B}\left(\epsilon_{\boldsymbol{p}}-\nu\right)+n_{F}\left(-\nu\right)\right]\left[n_{B}\left(\nu-\varepsilon_{\boldsymbol{q}}\right)+n_{F}\left(-\varepsilon_{\boldsymbol{q}}\right)\right]
\end{equation}

\begin{equation}
P_{\boldsymbol{p}\boldsymbol{q}}^{ep}=kT\frac{2\pi}{\hbar}\left(-\frac{\partial f_{\boldsymbol{p}}}{\partial\epsilon_{\boldsymbol{p}}}\right)\left|\alpha\left(\boldsymbol{p},\boldsymbol{q}\right)\right|^{2}\mathcal{B}\left(\boldsymbol{p}-\boldsymbol{q},\varepsilon_{\boldsymbol{p}}-\epsilon_{\boldsymbol{q}}\right)\left(n_{B}\left(\varepsilon_{\boldsymbol{p}}-\epsilon_{\boldsymbol{q}}\right)+f\left(-\epsilon_{\boldsymbol{q}}\right)\right)
\end{equation}

where $\chi''\left(\boldsymbol{p},\omega\right)$ is the spin susceptibility
and $\mathcal{B}\left(\boldsymbol{k},\nu\right)$ is the phonon spectral
function. We also reintroduced the momentum dependence in all the
couplings. In the case where the magnetic and itinerant electron layers
are perfectly aligned we have 
\begin{equation}
J_{K}\left(\boldsymbol{p},\boldsymbol{q}\right)=J_{K}\left[2\cos\left(p_{z}a_{z}\right)\right]\left[2\cos\left(q_{z}a_{z}\right)\right],
\end{equation}

\begin{equation}
\tilde{\alpha}\left(\boldsymbol{p},\boldsymbol{q}\right)=\tilde{\alpha}\left[2\cos\left(p_{z}a_{z}\right)\right]\left[2\cos\left(q_{z}a_{z}\right)\right],
\end{equation}

and

\begin{equation}
\alpha\left(\boldsymbol{p},\boldsymbol{q}\right)=\alpha\left[2\cos\left(p_{z}a_{z}\right)\right]\left[2\cos\left(q_{z}a_{z}\right)\right].
\end{equation}

For the c axis resistivity we also consider the dispersion 
\begin{equation}
\epsilon_{\boldsymbol{k}}=\varepsilon_{\boldsymbol{k}}-2t_{\perp}\cos\left(k_{z}a_{z}\right)
\end{equation}

where $\varepsilon_{\boldsymbol{k}}$ is the purely 2d dispersion
considered in the main text. Additonally we consider the local approximation
for 
\begin{equation}
\chi''\left(\nu\right)=\kappa\nu\frac{\theta\left(\alpha\pi J_{H}-\left|\nu\right|\right)}{J_{H}}
\end{equation}

with the sum rule constraint 
\begin{equation}
\int\frac{\chi''\left(\nu\right)}{\nu}d\nu=2\pi S\left(S+1\right).
\end{equation}

Focusing now on the case of optical phonons
\begin{equation}
\mathcal{B}\left(\boldsymbol{k},\nu\right)=\frac{\pi\hbar}{2M\omega_{0}}\left(\delta\left(\nu-\omega_{0}\right)-\delta\left(\nu+\omega_{0}\right)\right)
\end{equation}

we can then calculate the transport time for the coherent conductivity
channel. The integrals above can then be evaluated analytically and
the resistivity becomes

\begin{align}
\rho_{\perp} & =\frac{\hbar}{e^{2}}\frac{a_{xy}^{2}}{a_{z}}\left(\frac{1}{2\nu_{0}\left(2t_{\perp}\right)^{2}}\right)\left[2\pi\lambda_{K}J_{K}\mathcal{I}\left(\alpha\pi\beta J\right)+2\pi\frac{\lambda_{tr,ME}}{\beta}\mathcal{M}\left(\alpha\pi\beta J,\beta\omega_{0}\right)+2\pi\frac{\lambda_{tr,E}}{\beta}\mathcal{K}\left(\beta\omega_{0}\right)\right]
\end{align}

With where the functions above Explicitly we have the definitions:

\begin{align}
\mathcal{I}\left(y\right) & =-\frac{4\text{Li}_{2}\left(e^{-y}\right)}{y}-\frac{2y}{e^{y}-1}-4y+\frac{2\pi^{2}}{3y}+4\log\left(e^{y}-1\right),
\end{align}

\begin{equation}
\mathcal{K}\left(x\right)=\left[\frac{x}{2}\text{csch}\left(\frac{x}{2}\right)\right]^{2},
\end{equation}

\begin{equation}
\mathcal{M}\left(y,z\right)=\frac{z^{2}}{\cosh\left(z\right)-1}\mathcal{J}\left(y,z\right),
\end{equation}

and,

\begin{align}
\mathcal{J}\left(y,z\right) & =\frac{-2z\text{Li}_{2}\left(e^{-y}\right)+\left(2y+z\right)\text{Li}_{2}\left(e^{-y-z}\right)+\left(z-2y\right)\text{Li}_{2}\left(e^{z-y}\right)+2\text{Li}_{3}\left(e^{-y-z}\right)-2\text{Li}_{3}\left(e^{z-y}\right)}{2yz}\\
 & +\frac{z^{2}}{12y}-\frac{1}{2}\log\left(\left(e^{y}-e^{z}\right)\left(e^{y+z}-1\right)\right)+\frac{y\log\left(\sinh\left(\frac{y-z}{2}\right)\text{csch}\left(\frac{y+z}{2}\right)\right)}{2z}+\frac{y}{2}+\frac{\pi^{2}}{3y}+\log\left(e^{y}-1\right)+\frac{z}{2}.
\end{align}

These functions have the following high temperature limits
\begin{equation}
\lim_{T\rightarrow\infty}\mathcal{M}\left(\alpha\pi\beta J,\beta\omega_{0}\right)=\lim_{T\rightarrow\infty}\mathcal{K}\left(\beta\omega_{0}\right)=1
\end{equation}

and 
\begin{equation}
\lim_{T\rightarrow\infty}\mathcal{I}\left(\alpha\pi\beta J\right)=2.
\end{equation}

\end{document}
